% Appendix C — Notas teóricas

\chapter{Notas teóricas}
\label{AppendC}

Este apéndice reúne las dos derivaciones que sustentan el pipeline pero que se
omitieron del cuerpo por brevedad: el modelo de difusión que genera la imagen y
la solución cerrada del regresor Ridge que mapea fMRI al espacio CLIP.

\section{Modelos de difusión: derivación abreviada}

\paragraph{Proceso directo (forward).} Un DDPM \cite{ho2020ddpm} define una
cadena de Markov que agrega ruido gaussiano a la imagen $x_0$ en $T$ pasos según
una varianza $\{\beta_t\}_{t=1}^{T}$:
\begin{equation}
  q(x_t \mid x_{t-1}) = \mathcal{N}\!\left(x_t;\, \sqrt{1-\beta_t}\,x_{t-1},\, \beta_t \mathbf{I}\right).
\end{equation}
Con $\alpha_t = 1-\beta_t$ y $\bar\alpha_t = \prod_{s=1}^{t}\alpha_s$, la
marginal admite muestreo directo en un paso:
\begin{equation}
  x_t = \sqrt{\bar\alpha_t}\,x_0 + \sqrt{1-\bar\alpha_t}\,\epsilon,\qquad \epsilon \sim \mathcal{N}(\mathbf{0},\mathbf{I}).
\end{equation}

\paragraph{Proceso inverso (reverse).} Una red $\epsilon_\theta(x_t,t)$ aprende a
predecir el ruido; el entrenamiento minimiza el objetivo simplificado
\begin{equation}
  \mathcal{L}_\text{simple} = \mathbb{E}_{x_0,\epsilon,t}\!\left[\,\bigl\|\epsilon - \epsilon_\theta(x_t,t)\bigr\|^2\,\right].
\end{equation}
El muestreo parte de $x_T \sim \mathcal{N}(\mathbf{0},\mathbf{I})$ y aplica la
media posterior estimada paso a paso hasta $x_0$.

\paragraph{Guía sin clasificador (CFG).} Para condicionar la generación por un
vector $c$ (aquí el embedding CLIP predicho desde fMRI) se combina la predicción
condicional e incondicional con una escala de guía $w$
\cite{ho2022cfg}:
\begin{equation}
  \tilde\epsilon_\theta(x_t,t,c) = (1+w)\,\epsilon_\theta(x_t,t,c) - w\,\epsilon_\theta(x_t,t,\varnothing).
\end{equation}
En este trabajo $c$ es el \emph{image embedding} CLIP ViT-L/14 y el generador es
SD~2.1 unCLIP \cite{ramesh2022unclip,rombach2022ldm}; el valor
\emph{guidance\_scale}$=8.0$ (Ap.~\ref{AppendA}) fija el compromiso entre fidelidad
al vector del cerebro y coherencia visual.

\section{Derivación de la solución cerrada de Ridge}

El adapter lineal estima $W \in \mathbb{R}^{d\times p}$ que mapea $p$ vóxeles a
las $d=768$ dimensiones del espacio CLIP. Dadas $X \in \mathbb{R}^{n\times p}$
(betas) e $Y \in \mathbb{R}^{n\times d}$ (targets CLIP), la regresión Ridge
\cite{hoerl1970ridge} minimiza el error cuadrático con penalización $\ell_2$:
\begin{equation}
  \mathcal{L}(W) = \lVert XW - Y \rVert_F^2 + \lambda \lVert W \rVert_F^2.
\end{equation}
El objetivo es convexo y diferenciable; su gradiente es
\begin{equation}
  \nabla_W \mathcal{L} = 2X^\top(XW - Y) + 2\lambda W.
\end{equation}
Igualando a cero, $X^\top X W + \lambda W = X^\top Y$, y factorizando $W$ se
obtiene la solución cerrada
\begin{equation}
  W^\star = \left(X^\top X + \lambda \mathbf{I}\right)^{-1} X^\top Y.
\end{equation}
El término $\lambda\mathbf{I}$ garantiza que la matriz sea invertible incluso con
$p > n$ (más vóxeles que trials), régimen habitual en fMRI. El \emph{shrinkage}
que introduce $\lambda$ reduce la varianza a costa de sesgar la magnitud del
embedding hacia cero; por ello \texttt{visual\_evaluator.py} ofrece el modo de
normalización \texttt{ridge}, que reescala la norma antes de inyectar el vector a
SD~2.1 unCLIP. La variante \texttt{adapter\_ridge\_stoch.py} sustituye la
inversión exacta por descenso estocástico, habilitando penalizaciones no
cerradas y constituyendo el peldaño de contribución hacia los talleres
siguientes.
